\documentclass[12pt,a4paper]{article}
\usepackage[utf8]{inputenc}
\usepackage[spanish,es-noquoting]{babel}
\usepackage[T1]{fontenc}
\usepackage{lmodern}
\usepackage{geometry}
\geometry{top=2cm, bottom=2cm, left=2.5cm, right=2.5cm}
\usepackage{xcolor}
\usepackage{listings}
\usepackage{enumitem}
\usepackage{fancyhdr}
\usepackage{hyperref}
\usepackage{tikz}
\usetikzlibrary{arrows.meta, positioning, calc}

\definecolor{codegreen}{rgb}{0,0.6,0}
\definecolor{backcolour}{rgb}{0.95,0.95,0.92}

\lstset{
    language=C,
    backgroundcolor=\color{backcolour},
    commentstyle=\color{codegreen},
    keywordstyle=\color{blue},
    basicstyle=\ttfamily\small,
    breaklines=true,
    frame=single,
    numbers=left,
    tabsize=4,
    extendedchars=true,
    showstringspaces=false,
    inputencoding=utf8,
    literate={á}{{\'a}}1 {é}{{\'e}}1 {í}{{\'i}}1 {ó}{{\'o}}1 {ú}{{\'u}}1 {ñ}{{\~n}}1 {¡}{{\textexclamdown}}1 {¿}{{\textquestiondown}}1
}

\pagestyle{fancy}
\fancyhf{}
\lhead{Monitorías Sistemas Operativos 2026-I}
\rhead{Capítulo 3: Memoria en Hilos}
\cfoot{\thepage}

\title{\textbf{Guía de Aprendizaje: Modelo de Memoria en Hilos}}
\author{Malak Sanchez \\ \small Monitorías de Sistemas Operativos}
\date{\today}

\begin{document}

\maketitle

\section{Compartir vs Aislar}
A diferencia de los procesos (donde cada uno tiene su "isla" privada de memoria), los hilos viven en una casa compartida. Entender qué es privado y qué es público es vital para evitar errores desastrosos.

\section{Mapa de Memoria}
Todos los hilos de un mismo proceso comparten el direccionamiento virtual.

\begin{center}
\begin{tikzpicture}[font=\small, >=Stealth, 
    seg/.style={draw, thick, fill=gray!10, minimum width=3cm, minimum height=0.7cm, align=center},
    shared/.style={fill=green!15},
    priv/.style={fill=blue!15}]

    % Memory stack
    \node[seg, shared] (heap) at (0, 3) {Heap (Shared)};
    \node[seg, shared] (data) at (0, 2.4) {Global Data (Shared)};
    \node[seg, shared] (code) at (0, 1.8) {Code / Text (Shared)};

    % Gap
    \node at (0, 1.2) {...};

    % Individual Stacks
    \node[seg, priv] (s1) at (-2, 0) {Stack Thread 1 \\ \tiny (Local Vars)};
    \node[seg, priv] (s2) at (2, 0) {Stack Thread 2 \\ \tiny (Local Vars)};

    % Connections
    \draw[<->, dashed, gray] (heap) -- (s1);
    \draw[<->, dashed, gray] (heap) -- (s2);

\end{tikzpicture}
\end{center}

\section{¿Qué se comparte?}
\begin{itemize}
    \item \textbf{Variables Globales:} Si un hilo cambia una variable global, el otro verá el cambio inmediatamente.
    \item \textbf{Heap (\texttt{malloc}):} Si un hilo reserva memoria, el puntero es válido para todos los demás.
    \item \textbf{Descriptores de Archivo:} Si un hilo abre un archivo, otro puede leerlo usando el mismo FD.
\end{itemize}

\section{¿Qué es privado?}
\begin{itemize}
    \item \textbf{Stack:} Las variables declaradas dentro de la función del hilo son privadas.
    \item \textbf{Registros / PC:} Cada hilo está en su propia línea de código.
\end{itemize}

\section{Ejemplo: La Variable Problemática}
\begin{lstlisting}
#include <stdio.h>
#include <pthread.h>

int global_shared = 0; // Public

void* worker(void* arg) {
    global_shared++; // Modifying shared memory
    return NULL;
}

int main() {
    pthread_t t1, t2;
    pthread_create(&t1, NULL, worker, NULL);
    pthread_create(&t2, NULL, worker, NULL);
    pthread_join(t1, NULL);
    pthread_join(t2, NULL);
    printf("Final value: %d\n", global_shared);
    return 0;
}
\end{lstlisting}

\section{Ejercicios}
\begin{enumerate}
    \item Modifique el ejemplo anterior para que cada hilo incremente la variable 1,000,000 de veces. ¿El resultado es siempre 2,000,000? ¿Por qué?
    \item Intente pasar un puntero a una variable local del \texttt{main} a dos hilos diferentes. ¿Pueden ambos modificarla?
\end{enumerate}

\end{document}
